Esta interpretación la redactó Claude a partir de los registros de la corrida; no es del autor.

Luna pasó por dos fases. En el episodio~1 propuso sólo anotaciones: con la segunda, el juez generó 10 obligaciones y la automatización cerró 7. Desde el paso~3 del episodio~1 escribió lemas auxiliares, y ningún intento volvió a llegar a las obligaciones: en los 10 intentos restantes Isabelle rechazó algún lema. Nueve de esos diez intentos contienen el mismo enunciado, con dos nombres y dos formas equivalentes: para $x > 1$ real, $\lceil \log_2 (x/2) \rceil + 1 = \lceil \log_2 x \rceil$ (con los techos pasados a naturales). El enunciado es cierto, y es el mismo tipo de hecho, el techo de $\log_2$ al partir el intervalo a la mitad, que resuelve el lema \texttt{var\_dec} que los autores del artículo demostraron a mano con una línea \texttt{smt} de 18 hechos.

Los errores de Isabelle muestran fallas de conocimiento del sistema más que de matemática. En tres intentos Luna usó métodos de otros asistentes de prueba (\texttt{omega}; \texttt{field\_simp}, cuyo análogo en Isabelle es \texttt{field\_simps}) y en dos citó hechos que no existen (\texttt{divide\_less\_iff}; \texttt{log\_div}, cuando el nombre correcto es \texttt{log\_divide}). Los otros cinco son demostraciones que no cierran (\emph{Failed to finish proof}).

La memoria hizo su trabajo: la biblioteca terminó con 4 lemas que Luna demostró e Isabelle verificó, y cada intento los recibió. Pero el diseño del juez convierte un lema fallido en un intento perdido: si un lema del bloque falla, el juez no evalúa las obligaciones. Por eso, desde el episodio~2, cada intento se jugó entero a ese único lema. Evaluar las obligaciones con los lemas que sí se verificaron es un diseño alternativo razonable; no se cambió a mitad de la comparación porque Astra corrió con el mismo juez.
